\chapter{Criterio de familia operativa}

\section{Familia operativa}
\begin{definition}
Se llama familia operativa al conjunto disciplinado de configuraciones que comparten herencia canónica suficiente, compatibilidad de tratamiento y una forma legible de articulación dentro del canon intermedio.
\end{definition}

\begin{observation}
Una familia operativa no es todavía una clasificación exhaustiva del sistema. Su función es reunir configuraciones tratables bajo una misma gramática sin colapsar sus diferencias internas.
\end{observation}

\section{Condición mínima de familia}
\begin{definition}
Existe familia operativa cuando concurren simultáneamente:
\begin{enumerate}
  \item dependencia respecto de una base canónica común o compatible;
  \item comparabilidad suficiente entre sus configuraciones;
  \item distinción operativa conservada entre sus miembros;
  \item articulación jerárquica o relacional no contradictoria;
  \item función común reconocible dentro del canon intermedio.
\end{enumerate}
\end{definition}

\begin{observation}
La familia no exige identidad fuerte entre sus miembros. Exige solo una compatibilidad estructural suficiente para tratarlos sin arbitrariedad.
\end{observation}

\section{Familia y configuración}
\begin{definition}
Toda familia operativa se compone de configuraciones que ya poseen estatus tratable dentro del Tomo II.
\end{definition}

\begin{observation}
Esta cláusula impide formar familias a partir de piezas todavía indeterminadas o sin umbral operativo suficiente.
\end{observation}

\section{Familia y jerarquía}
\begin{definition}
Una familia operativa puede contener configuraciones de distinto nivel siempre que la relación entre esos niveles sea traducible y disciplinadamente articulable.
\end{definition}

\begin{observation}
La diversidad jerárquica dentro de una familia no es un defecto por sí misma. Se vuelve problemática solo cuando la heterogeneidad rompe comparabilidad o articulación.
\end{observation}

\section{Condición mínima de fallo}
\begin{definition}
Se considera fallida una familia operativa cuando sus miembros dejan de compartir herencia suficiente, función común reconocible o posibilidad de relación disciplinada estable.
\end{definition}

\begin{observation}
Esta condición protege al canon intermedio frente a agrupaciones nominales o puramente decorativas.
\end{observation}

\section{Cláusula de alcance}
\begin{observation}
El criterio de familia operativa fija una forma mínima de agrupación disciplinada. No constituye todavía una taxonomía exhaustiva de todas las familias posibles del sistema.
\end{observation}

